\documentclass[11pt]{article}
\usepackage[letterpaper,margin=1in]{geometry}
\usepackage[T1]{fontenc}
\usepackage{lmodern,microtype,amsmath,amssymb,amsthm}
\usepackage{xcolor}
\usepackage[colorlinks=true,linkcolor=black,citecolor=blue!45!black,urlcolor=blue!45!black]{hyperref}
\usepackage{xurl}
\newtheorem{theorem}{Theorem}[section]
\newtheorem{lemma}[theorem]{Lemma}
\newtheorem{corollary}[theorem]{Corollary}
\theoremstyle{remark}
\newtheorem{remark}[theorem]{Remark}
\newcommand{\DTIME}{\mathsf{DTIME}}
\newcommand{\DSPACE}{\mathsf{DSPACE}}
\newcommand{\NTIME}{\mathsf{NTIME}}
\newcommand{\SAT}{\textup{\textsc{CNF-SAT}}}
\hypersetup{pdftitle={Simulating time T with a space-s oracle in space O(s(T)+T/log T)},
  pdfauthor={Charlie Cruz},
  pdfkeywords={time-space simulation, oracle machines, relativization, space complexity, multitape Turing machines, CNF-SAT}}
\title{Simulating Time $T$ with a Space-$s$ Oracle\texorpdfstring{\\}{ }in Space $O(s(T)+T/\log T)$}
\author{Charlie Cruz}
\date{November 24, 2026}
\begin{document}
\maketitle
\begin{abstract}
Williams showed that multitape time $t$ can be simulated in space
$O(\sqrt{t\log t})$, and that the simulation survives an oracle provided
every query is short. We treat long queries. An oracle machine makes
\emph{fresh queries} if it writes each query from scratch, one symbol per
step, on a query tape that is erased after the answer arrives. We show
that a deterministic machine running in time $T$, with adaptive fresh
queries to an oracle $A$ decidable in space $s$, can be simulated without
any oracle in space $O(s(T)+T/\log T)$. The simulation stores neither the
machine's tapes nor its queries. It keeps a table of the answers to all
short queries and a transcript of the answers to the long ones. Whenever
the algorithm for $A$ asks for a bit of a long query, the simulation
regenerates that bit by replaying the whole computation inside
Williams's simulation. For example, \SAT{} with binary variable names is
in space $O(N/\log N)$, so linear space is not contained in time
$o(n\log n)$ with fresh queries to a \SAT{} oracle.
\end{abstract}

\section{Introduction}

\subsection{Background}
Hopcroft, Paul and Valiant showed that
$\DTIME[t]\subseteq\DSPACE[t/\log t]$ \cite{HPV}. Nearly fifty years
later, Williams improved this to
$\DTIME[t]\subseteq\DSPACE[\sqrt{t\log t}]$ \cite{Wil25}, building on the
tree evaluation algorithm of Cook and Mertz \cite{CM}. Both results are
for multitape Turing machines.

We ask what happens when the time-bounded machine has an oracle. Oracle
machines model algorithms that call subroutines whose internals we do
not want to examine, so it is natural to ask whether the square-root
simulation is robust to such calls. Williams observed that his
simulation relativizes in a restricted sense: if every query has length
at most $\sqrt{t\log t}$, then each query fits inside one block of the
simulation and can be passed along \cite[Theorem~5.1]{Wil25}. A machine
running in time $T$ can also write queries of length $\Theta(T)$, and it
can ask many of them, each depending on all the previous answers. A long
query depends on an entire region of the tape at once, which destroys
the locality that drives every known simulation. The obvious simulation
stores the whole configuration, which costs space $\Theta(T)$.

If the simulator keeps the oracle and is only required to use small
space, nothing can be gained. A small-space machine cannot even write
the queries that a time-$T$ machine can, and there are oracles relative
to which time and space coincide \cite{HCCRR,CW04}. We therefore ask a
different question. Suppose the oracle language $A$ is itself decidable
in space $s$. Can we remove the oracle and still use space less than
$T$?

\subsection{Results}
The answer is yes, provided each query has to be written out. Say that
an oracle machine writes \emph{fresh queries} if it writes each query
from scratch on a query tape that is erased after the answer arrives,
paying one step per symbol. Some convention of this kind is necessary;
see Remark~\ref{rem:fresh}.

\begin{theorem}\label{thm:main}
Let $A\subseteq\{0,1\}^*$ be decidable in space $s(m)$, where
$s(m)\ge\log m$ is nondecreasing and $s(2m)=O(s(m))$. Then for every
$T(n)\ge n$, every language decided by a deterministic oracle machine
with fresh queries to $A$ in time $T$ is in
\[
\DSPACE\bigl[O\bigl(s(T)+T/\log T\bigr)\bigr].
\]
\end{theorem}

The two terms have a direct interpretation. The term $s(T)$ is the cost
of answering a single query of length up to $T$. The term $T/\log T$ is
the number of \emph{long} queries the machine can afford to write. Their
answers are the only information in the computation that we cannot
regenerate from scratch, so we write them down. Everything else, namely
the machine's work tapes, its short queries and the contents of its long
queries, is either tabulated or recomputed on demand within space
$O(T^{2/3}\log T)$.

Theorem~\ref{thm:main} uses nothing about $A$ except its space
complexity. For an explicitly encoded CNF formula of length $N$, there
are only $O(N/\log N)$ variables, so \SAT{} is in space $O(N/\log N)$
(Lemma~\ref{lem:sat}). This gives the following.

\begin{corollary}\label{cor:sat}
For every $t(n)\ge n$ with $t(n)=o(n\log n)$, linear space is not
contained in deterministic time $t$ with fresh queries to a \SAT{} oracle:
\[
\DTIME^{\SAT}[t]\subsetneq\DSPACE[n].
\]
\end{corollary}

Corollary~\ref{cor:sat} is a linear-time analogue of
$\mathsf{P}^{\mathsf{NP}}\neq\mathsf{PSPACE}$: a machine that can solve
any satisfiability instance in one step still cannot decide all of
linear space in slightly superlinear time. Its scope is limited. The
NP-hardness of \SAT{} plays no role, since only its space complexity is
used. The $O(N/\log N)$ space bound comes from writing variable names in
binary, and the result is sensitive to that encoding. The corollary says
nothing about $\mathsf{P}^{\mathsf{NP}}$ versus $\mathsf{PSPACE}$, because
separations do not pad upward. Still, we do not see how to prove
Corollary~\ref{cor:sat} without a result like Theorem~\ref{thm:main}.
Diagonalizing over oracle machines in linear space requires simulating
them in sublinear space, and that is what was missing.

The same argument applies to any oracle that is decidable in slightly
sublinear space. For example:

\begin{corollary}\label{cor:ntime}
Suppose $\NTIME[m]\subseteq\DSPACE[O(m/\sqrt{\log m})]$. Then for every
$A\in\NTIME[n]$ and every $t(n)\ge n$ with $t(n)=o(n\sqrt{\log n})$,
\[
\DTIME^{A}[t]\subsetneq\DSPACE[n].
\]
\end{corollary}

\subsection{Techniques}
Long queries cause difficulty only because we cannot afford to store
them, so we do not store them. When the algorithm for $A$ asks for bit
$i$ of the current query, we rerun the whole computation from the
beginning, feeding in the stored answers, and stop when bit $i$ of that
query is written. The rerun is a long deterministic computation. It is
slightly longer than the original, because short queries are answered by
scanning a table, but it uses no oracle, so Williams's simulation runs it
in roughly the square root of its time. The older bound of \cite{HPV}
would not suffice here. Because short queries are answered by scanning a
table, the rerun takes about $T^{4/3}$ steps, so \cite{HPV} would need
space about $T^{4/3}/\log T$ to regenerate a single bit, which exceeds
$T$. The argument works because Williams's simulation regenerates a bit in
space polynomially smaller than $T$.

\subsection{Status of imported results}
The proof of Theorem~\ref{thm:main} uses Williams's simulation
(Theorem~\ref{thm:williams}) as a black box, and Corollary~\ref{cor:sat}
also uses the space hierarchy theorem \cite{SHL}. Corollary~\ref{cor:ntime}
is conditional on its hypothesis. That hypothesis is claimed in a recent
preprint of Ben Daniel \cite{BD26}. It also follows from the alternating
simulation in our note \cite{AS}, which uses random access to the input,
combined with the classical fact that random-access alternating time $t$
is in deterministic space $O(t)$ \cite{CKS}.

\section{Preliminaries}
Machines are multitape Turing machines over the binary alphabet, apart
from delimiters. An oracle machine also has a write-only query tape. To
ask a query, it writes a string $z\in\{0,1\}^*$ on that tape from the left
end, one cell per step, and enters a query state. In one step it then
receives $A(z)$, and the query tape is erased. So a query of length $m$
costs at least $m$ steps.

We use the following standard fact about composing space-bounded
computations. Suppose a space-$f$ machine reads its input $y$ through a
read-only input tape, and each symbol $y_j$ can be computed from $j$ in
space $g$. Then the machine can be run on $y$ in space
$O(f+g+\log|y|)$ without ever writing $y$ down: keep the position of
the input head as a counter, and recompute $y_j$ whenever it is read.

\begin{theorem}[\cite{Wil25}]\label{thm:williams}
For every $t(n)\ge n$, $\DTIME[t]\subseteq\DSPACE[O(\sqrt{t\log t})]$.
\end{theorem}

\section{The simulation}
\begin{proof}[Proof of Theorem~\ref{thm:main}]
Let $M$ be an oracle machine with fresh queries running in time $T(n)$,
and fix an input $x$ of length $n$. We first suppose that we are given a
power of two $U$ with $n\le U$ and $T(n)\le U$, and we remove this
assumption at the end.

\emph{Short and long queries.}
Put $h=\lfloor(\log U)/3\rfloor$. A query is \emph{short} if its length is
at most $h$, and \emph{long} otherwise. Writing a long query costs more
than $h$ steps, so $M$ makes at most $U/(h+1)=O(U/\log U)$ long queries.

\emph{What we store.}
We store a table $\tau$ of the bits $A(z)$ for all $z\in\{0,1\}^{\le h}$,
in lexicographic order. It has at most $2^{h+1}\le2U^{1/3}$ bits, and we
fill it in by running the space-$s$ algorithm for $A$ on each short
string, reusing the same workspace. We also store a transcript $\beta$
of the answers to the first $|\beta|$ long queries of $M^A$ on $x$. It
starts empty and never exceeds $O(U/\log U)$ bits.

\emph{The replay machine.}
Let $R$ be the following deterministic machine without an oracle. Its
input has the form $y=\langle x,\tau,\beta,i,1^U\rangle$. It first checks
that $\tau$ has length $2^{h+1}-1$ and that $|\beta|\le U$, and it rejects
otherwise, so that its running time is bounded on all inputs. It simulates $M$
on $x$, using the padding $1^U$ as a clock, and stops after $U$ steps of
$M$. It answers each short query by scanning $\tau$. It answers the
$\ell$-th long query with $\beta_\ell$, keeping a head on $\beta$ that
moves right once per long query. When $M$ halts, times out, or completes
its $(|\beta|+1)$-st long query $z$, the machine $R$ stops. It then reports
which of these happened, together with $M$'s answer, or $|z|$ in binary,
or the bit $z_i$. We regard $R$ as deciding the language of pairs
$(y,j)$ for which bit $j$ of this report is $1$.

As long as $\beta$ holds true answers, $R$ follows the actual
computation of $M^A$ on $x$. Each short-query lookup costs
$O(2^hh)$ steps, and there are at most $U$ queries, so $R$ halts within
$O(U^{4/3}\log U)$ steps. Its input length is $\Theta(U)$, because of the
padding. So $R$ runs in time $O(|y|^{4/3}\log|y|)$, and by
Theorem~\ref{thm:williams} it can be simulated in space
$O(|y|^{2/3}\log|y|)=O(U^{2/3}\log U)$. We never write $y$ down. Its
symbols come from the true input, from the stored $\tau$ and $\beta$, from
a counter holding $i$, and from the padding, which is generated on the
fly.

\emph{The main loop.}
Repeat the following.
\begin{enumerate}
\item Read the report of $R$ by simulating it. If $M$ halted, output its
answer and stop. If $M$ timed out, abandon this value of $U$.
\item Otherwise $M$ has just written a long query $z$ of known length
$m\le U$. Run the space-$s$ algorithm for $A$ on $z$, which exists only
virtually. Each time that algorithm reads position $i$ of $z$, simulate
$R$ with index $i$ to compute $z_i$.
\item Append $A(z)$ to $\beta$.
\end{enumerate}
By induction, $\beta$ always holds true answers, so the loop follows
$M^A$ on $x$ exactly. It runs for at most $O(U/\log U)$ iterations.

\emph{Space.}
The table takes $O(U^{1/3})$ space and the transcript $O(U/\log U)$. The
algorithm for $A$ takes $O(s(U))$, plus an index of $O(\log U)$ bits.
The simulation of $R$ takes $O(U^{2/3}\log U)$ and is reused across
calls. The total is $O(s(U)+U/\log U)$.

\emph{Removing the clock.}
Run the procedure for $U=2^{\lceil\log n\rceil},2^{\lceil\log n\rceil+1},\dots$
until some run does not time out. Every run uses true oracle answers
only, so the first run to finish gives the right output. A run with
$U\ge T(n)$ never times out, so the last $U$ is at most the larger of
$2^{\lceil\log n\rceil}<2n$ and the least power of two that is at least
$T(n)$. Either way it is less than $2T(n)$. Space is reused between runs,
so the total is $O(s(2T)+T/\log T)=O(s(T)+T/\log T)$.
\end{proof}

\begin{remark}\label{rem:fresh}
The fresh-query convention is used in exactly one place, to bound the
number of long queries by $O(U/\log U)$. If a machine could edit a long
query in place and ask it again, then $T$ steps would buy $\Omega(T)$
long queries, and the transcript alone would take linear space. Any
convention that charges at least $h$ steps per long query works as
well.
\end{remark}

\section{Applications}
\begin{lemma}\label{lem:sat}
Encode a CNF formula as a list of clauses, each literal given by a sign
bit and a binary variable identifier, with delimiters coded over the
binary alphabet. Then \SAT{} is in $\DSPACE[O(N/\log N)]$ on inputs of
length $N$.
\end{lemma}
\begin{proof}
Let $v$ be the number of distinct identifiers. Fewer than $v/2$ binary
strings have length less than $\lfloor\log v\rfloor-1$, so at least $v/2$
identifiers have length at least $\log v-2$. Hence $N\ge(v/2)(\log v-2)$,
and $v=O(N/\log N)$.

Run through all assignments $a\in\{0,1\}^v$ with a $v$-bit counter. The
number $v$, and the rank of any identifier among the distinct
identifiers, can be computed in space $O(\log N)$ by comparing
occurrences pairwise and counting first occurrences. For each assignment,
evaluate every clause, looking up each literal in $a$ by its rank. The
space is $v+O(\log N)$.
\end{proof}

\begin{proof}[Proof of Corollary~\ref{cor:sat}]
By Lemma~\ref{lem:sat} and Theorem~\ref{thm:main} with $s(m)=m/\log m$,
a machine running in time $t$ can be simulated in space
$O(t/\log t)\le O(t/\log n)=o(n)$. By the space hierarchy theorem
\cite{SHL}, some language in $\DSPACE[n]$ is not in $\DSPACE[o(n)]$.
\end{proof}

\begin{proof}[Proof of Corollary~\ref{cor:ntime}]
Use Theorem~\ref{thm:main} with $s(m)=m/\sqrt{\log m}$. The space is
$O(t/\sqrt{\log t})=o(n)$ for $t=o(n\sqrt{\log n})$, and we conclude as
before.
\end{proof}

\section{Discussion}
\paragraph{Related work.}
Theorem~5.1 of \cite{Wil25} covers the case where all queries
have length at most $\sqrt{t\log t}$, with the simulator keeping access to
the oracle. Chen and Wang make the same observation, under the name
localization \cite{CW26}. Tretkoff studied length-restricted
relativizations of earlier simulations inside linear space \cite{Tre86}.
We have not been able to check whether that paper treats long queries
against a space-bounded oracle. A search of the follow-ups to
\cite{Wil25} turned up no treatment of oracle machines with long queries.

\paragraph{Open problems.}
\begin{itemize}
\item \emph{Can the transcript be compressed?} The $T/\log T$ term is the
only place where adaptivity costs anything. Is
$\DTIME^{\SAT}[t]\subseteq\DSPACE[o(t/\log t)]$? This would push
Corollary~\ref{cor:sat} beyond time $n\log n$. A lower bound for this
kind of simulation would also be interesting, since the answer bits
look incompressible but we see no proof that they are.
\item \emph{Nondeterministic outer machines.} For $\NTIME^A$ the
simulator must check guessed answers. When $A$ has certificates for both
answers, for instance $A\in\NTIME[n]\cap\mathsf{coNTIME}[n]$, the standard
argument that multiplies gap functions gives a linear-time reduction to
counting. The general case is the second level of the hierarchy, and it
seems to require new ideas.
\item \emph{Other models.} Williams's simulation has recently been
extended beyond multitape machines \cite{SW26}. Does
Theorem~\ref{thm:main} extend in the same way?
\end{itemize}

\begin{thebibliography}{99}
\bibitem{AS}
C.~Cruz. Alternating speedups for nondeterministic computations. Note,
November 2026, Theorem~1.1.

\bibitem{BD26}
S.~Ben Daniel. Randomized, quantum and postselected multitape time in
deterministic space $t/\sqrt{\log t}$. arXiv:2611.01884, 2026.

\bibitem{CW04}
J.-Y. Cai and O.~Watanabe. On proving circuit lower bounds against the
polynomial-time hierarchy. \emph{SIAM Journal on Computing},
33(4):984--1009, 2004.

\bibitem{CKS}
A.~K. Chandra, D.~C. Kozen, and L.~J. Stockmeyer. Alternation.
\emph{Journal of the ACM}, 28(1):114--133, 1981.

\bibitem{CW26}
L.~Chen and Y.~Wang. Separating RAM and multitape Turing machines with
short random oracles. ECCC TR26-006, 2026.

\bibitem{CM}
J.~Cook and I.~Mertz. Tree evaluation is in space
$O(\log n\cdot\log\log n)$. In \emph{Proceedings of STOC 2024}, 2024.

\bibitem{HCCRR}
J.~Hartmanis, R.~Chang, S.~Chari, D.~Ranjan, and P.~Rohatgi.
Relativization: a revisionistic retrospective. In \emph{Current Trends in
Theoretical Computer Science}, World Scientific, 1993.

\bibitem{HPV}
J.~Hopcroft, W.~Paul, and L.~Valiant. On time versus space.
\emph{Journal of the ACM}, 24(2):332--337, 1977.

\bibitem{SW26}
D.~Sibgatullin and R.~R. Williams. Space-efficient simulations beyond
multitape Turing machines. ECCC TR26-223, 2026.

\bibitem{SHL}
R.~E. Stearns, J.~Hartmanis, and P.~M. Lewis~II. Hierarchies of memory
limited computations. In \emph{6th Annual Symposium on Switching Circuit
Theory and Logical Design}, 1965.

\bibitem{Tre86}
C.~Tretkoff. Bounded oracles and complexity classes inside linear space.
In \emph{Structure in Complexity Theory}, LNCS 223, pages 347--361, 1986.

\bibitem{Wil25}
R.~R. Williams. Simulating time with square-root space. In
\emph{Proceedings of STOC 2025}, 2025. arXiv:2502.17779.
\end{thebibliography}
\end{document}
